% =====================================================================
%  BÖLÜM 5: Önemli Dağılım Aileleri
%  Kaynak: Feller Vol.I-II; Hogg & Craig; Casella & Berger; Klenke
% =====================================================================
\chapter{Önemli Dağılım Aileleri}
\label{chp:dagilim_aileleri}

\bolumalinti{%
  Doğa, birkaç temel dağılımın sonsuz çeşitlemesinden ibarettir.%
}{William Feller, \textit{An Introduction to Probability Theory and Its Applications}}

% ---- Kazanımlar (TYMM) ----
\begin{kazanimlar}
Bu bölümün sonunda öğrenci:
\begin{itemize}[leftmargin=1.8em]
  \item Temel kesikli dağılımları (Bernoulli, Binom, Poisson, Geometrik,
        Negatif Binom, Hipergeometrik, Çok Terimli) tanımlayabilir;
        beklenti, varyans ve MÜF'lerini hesaplayabilir.
  \item Temel sürekli dağılımları (Uniform, Normal, Üstel, Gama, Beta,
        Cauchy, Lognormal, $\chi^2$, $t$, $F$) tanımlayabilir; türetme ilişkilerini açıklayabilir.
  \item Üstel aile kavramını tanımlayabilir; doğal parametre ve yeterli istatistik
        kavramlarıyla bağlantısını kurabilir.
  \item Yer-ölçek ailesi kavramını açıklayabilir; standartlaştırma işlemini uygulayabilir.
  \item Dağılımlar arasındaki türetme ve limit ilişkilerini (Poisson limiti,
        normal yaklaşım, $\chi^2$ toplamı) ispatlayabilir.
\end{itemize}
\end{kazanimlar}

% ---- Motivasyon ----
\begin{motivasyon}
Bölüm~\ref{chp:rasgele_degiskenler} ve~\ref{chp:beklenti}'de dağılım fonksiyonu
ve beklenti soyut düzeyde incelendi. Şimdi pratikte en sık karşılaşılan dağılım
ailelerini somutlaştırıyoruz. Bu dağılımlar teorik güzelliklerinin yanı sıra
istatistiksel modellemenin (Kısım~\ref{part:istatistik}) temel yapı taşlarıdır.

Her dağılım için standart bilgi şeması: PMF/PDF — parametreler — $E[X]$, $\mathrm{Var}(X)$ —
MÜF/KF — özellikler — tipik kullanım alanı.

\medskip
\textbf{Köprü.} Bölüm~\ref{chp:beklenti}'deki MÜF ve KF araçları burada her dağılım
için hesaplanacak; Bölüm~\ref{chp:buyuk_sayilar} ve~\ref{chp:merkezi_limit}'te
bu dağılımların asimptotik davranışı incelenecektir.
\end{motivasyon}

% =====================================================================
\section{Kesikli Dağılımlar}
\label{sec:kesikli_dagilimlar}
% =====================================================================

\subsection{Bernoulli ve Binom Dağılımı}
\label{subsec:bernoulli_binom}

\begin{tanim}[Bernoulli Dağılımı]
\label{def:bernoulli}
$0 < p < 1$ olsun. $X \sim \mathrm{Bernoulli}(p)$:
\[
  P(X = 1) = p, \quad P(X = 0) = 1 - p =: q.
\]
\textbf{Beklenti:} $E[X] = p$. \quad
\textbf{Varyans:} $\mathrm{Var}(X) = pq$. \quad
\textbf{MÜF:} $M_X(t) = q + pe^t$.
\end{tanim}

\begin{tanim}[Binom Dağılımı]
\label{def:binom}
$X \sim \Binom(n, p)$: $n$ bağımsız $\mathrm{Bernoulli}(p)$ denemesinde başarı sayısı.
\[
  P(X = k) = \binom{n}{k} p^k q^{n-k}, \quad k = 0, 1, \ldots, n.
\]
\textbf{Beklenti:} $E[X] = np$. \quad
\textbf{Varyans:} $\mathrm{Var}(X) = npq$. \quad
\textbf{MÜF:} $M_X(t) = (q + pe^t)^n$.
\end{tanim}

\begin{proof}[MÜF İspat]
$X = X_1 + \cdots + X_n$ ($X_i$ bağımsız Bernoulli).
$M_X(t) = \prod_{i=1}^n M_{X_i}(t) = (q+pe^t)^n$.
\end{proof}

\begin{teorem}[Binom Yeniden Üretkenliği]
\label{thm:binom_yeniden}
$X \sim \Binom(n, p)$, $Y \sim \Binom(m, p)$ bağımsız $\Rightarrow$
$X + Y \sim \Binom(n+m, p)$.
\end{teorem}

\begin{proof}
$M_{X+Y}(t) = (q+pe^t)^n(q+pe^t)^m = (q+pe^t)^{n+m}$.
\end{proof}

\begin{onerme}[Binom Mod ve Medyan]
$\lfloor (n+1)p \rfloor$ binom dağılımının modudur (birden fazla mod olabilir).
$n = 10$, $p = 0.3$: mod $= \lfloor 11 \cdot 0.3 \rfloor = 3$.
\end{onerme}

\subsection{Poisson Dağılımı}
\label{subsec:poisson}

\begin{tanim}[Poisson Dağılımı]
\label{def:poisson}
$\lambda > 0$ olsun. $X \sim \Pois(\lambda)$:
\[
  P(X = k) = \frac{\lambda^k e^{-\lambda}}{k!}, \quad k = 0, 1, 2, \ldots
\]
\textbf{Beklenti:} $E[X] = \lambda$. \quad
\textbf{Varyans:} $\mathrm{Var}(X) = \lambda$. \quad
\textbf{MÜF:} $M_X(t) = e^{\lambda(e^t - 1)}$.
\end{tanim}

\begin{proof}[Beklenti ve MÜF]
$E[X] = \sum_{k=0}^\infty k \frac{\lambda^k e^{-\lambda}}{k!}
= \lambda e^{-\lambda} \sum_{k=1}^\infty \frac{\lambda^{k-1}}{(k-1)!} = \lambda$.
$M_X(t) = \sum_{k=0}^\infty e^{tk} \frac{\lambda^k e^{-\lambda}}{k!}
= e^{-\lambda} \sum_{k=0}^\infty \frac{(\lambda e^t)^k}{k!} = e^{-\lambda} e^{\lambda e^t} = e^{\lambda(e^t-1)}$.
\end{proof}

\begin{teorem}[Binom'dan Poisson Limiti]
\label{thm:binom_poisson}
$n \to \infty$, $p = p_n \to 0$, $np_n \to \lambda > 0$ iken
$\Binom(n, p_n) \xrightarrow{d} \Pois(\lambda)$.
\end{teorem}

\begin{proof}
$P(X_n = k) = \binom{n}{k} p_n^k (1-p_n)^{n-k}$.
$np_n \to \lambda$, yani $p_n = \lambda_n/n$ ($\lambda_n = np_n \to \lambda$).
\begin{align*}
  \binom{n}{k} p_n^k (1-p_n)^{n-k}
  &= \frac{n(n-1)\cdots(n-k+1)}{k!} \cdot \frac{\lambda_n^k}{n^k} \cdot \left(1 - \frac{\lambda_n}{n}\right)^{n-k} \\
  &\to \frac{1}{k!} \cdot \lambda^k \cdot e^{-\lambda} = \frac{\lambda^k e^{-\lambda}}{k!}.
\end{align*}
(Her $k$ için: ilk çarpan $\to 1$, $(1-\lambda_n/n)^n \to e^{-\lambda}$, $(1-\lambda_n/n)^{-k} \to 1$.)
\end{proof}

\begin{onerme}[Poisson Yeniden Üretkenliği]
$X \sim \Pois(\lambda)$, $Y \sim \Pois(\mu)$ bağımsız $\Rightarrow$ $X+Y \sim \Pois(\lambda+\mu)$.
\end{onerme}

\begin{proof}
$M_{X+Y}(t) = e^{\lambda(e^t-1)} \cdot e^{\mu(e^t-1)} = e^{(\lambda+\mu)(e^t-1)}$.
\end{proof}

\begin{ornek}[Poisson Süreci]
Birim zamanda ortalama $\lambda$ müşteri gelen bir bankada,
$k$ müşteri gelme olasılığı $P(X=k) = e^{-\lambda}\lambda^k/k!$.
$\lambda = 3$: $P(X = 0) = e^{-3} \approx 0.050$, $P(X \leq 2) \approx 0.423$.
\end{ornek}

\subsection{Geometrik ve Negatif Binom}
\label{subsec:geometrik}

\begin{tanim}[Geometrik Dağılım]
\label{def:geometrik}
$X \sim \mathrm{Geom}(p)$: ilk başarıya kadar gereken deneme sayısı.
\[
  P(X = k) = q^{k-1} p, \quad k = 1, 2, 3, \ldots
\]
\textbf{Beklenti:} $E[X] = 1/p$. \quad \textbf{Varyans:} $\mathrm{Var}(X) = q/p^2$.
\textbf{MÜF:} $M_X(t) = pe^t / (1 - qe^t)$, $t < -\ln q$.
\end{tanim}

\begin{onerme}[Bellek Yoksunluğu]
\label{prop:bellek_yoksunlugu}
$X \sim \mathrm{Geom}(p)$ için her $m, n \geq 1$'de
$P(X > m+n \mid X > m) = P(X > n)$.
Kesikli dağılımlar arasında bellek yoksunluğu (memorylessness) özelliğine sahip
tek dağılım geometrik dağılımdır.
\end{onerme}

\begin{proof}
$P(X > k) = q^k$.
$P(X > m+n \mid X > m) = \frac{P(X > m+n)}{P(X > m)} = \frac{q^{m+n}}{q^m} = q^n = P(X > n)$.
Tersine: bu özelliği sağlayan tek kesikli dağılımın geometrik olduğu,
$P(X > k) = P(X > 1)^k$ çarpım özelliğiyle gösterilir.
\end{proof}

\begin{tanim}[Negatif Binom Dağılımı]
\label{def:neg_binom}
$X \sim \mathrm{NegBin}(r, p)$: $r$'inci başarıya kadar gereken deneme sayısı.
\[
  P(X = k) = \binom{k-1}{r-1} p^r q^{k-r}, \quad k = r, r+1, \ldots
\]
\textbf{Beklenti:} $E[X] = r/p$. \quad
\textbf{Varyans:} $\mathrm{Var}(X) = rq/p^2$.\\
$\mathrm{NegBin}(1, p) = \mathrm{Geom}(p)$.
\end{tanim}

\subsection{Hipergeometrik Dağılım}
\label{subsec:hiper}

\begin{tanim}[Hipergeometrik Dağılım]
\label{def:hipergeometrik}
$N$ elemanlı popülasyonda $K$ "başarılı" eleman var; $n$ elemanlık örneklem çekilsin
(yerine koymadan). $X$: örneklemdeki başarılı sayısı.
\[
  P(X = k) = \frac{\binom{K}{k}\binom{N-K}{n-k}}{\binom{N}{n}},
  \quad \max(0, n+K-N) \leq k \leq \min(n, K).
\]
\textbf{Beklenti:} $E[X] = nK/N$. \quad
\textbf{Varyans:} $\mathrm{Var}(X) = n \frac{K}{N}\frac{N-K}{N}\frac{N-n}{N-1}$.
\end{tanim}

\begin{onerme}[Binom Yaklaşımı]
$N \to \infty$, $K/N \to p$ iken $\mathrm{Hiper}(N, K, n) \to \Binom(n, p)$.
Yani büyük popülasyonda yerine koymadan örneklem ≈ yerine koyarak örneklem.
\end{onerme}

\subsection{Çok Terimli Dağılım}
\label{subsec:multinomial}

\begin{tanim}[Çok Terimli Dağılım]
\label{def:multinomial}
$n$ bağımsız deneme; her deneyde $k$ kategoriden biri $p_1, \ldots, p_k$ olasılıklarla seçilsin
($\sum p_i = 1$). $X_i$: $i$'inci kategorinin deneme sayısı.
$(X_1, \ldots, X_k) \sim \mathrm{Multinomial}(n; p_1, \ldots, p_k)$:
\[
  P(X_1 = n_1, \ldots, X_k = n_k) = \frac{n!}{n_1! \cdots n_k!} p_1^{n_1} \cdots p_k^{n_k},
  \quad \sum n_i = n.
\]
\textbf{Beklenti:} $E[X_i] = np_i$. \quad
\textbf{Varyans:} $\mathrm{Var}(X_i) = np_i(1-p_i)$. \quad
\textbf{Kovaryans:} $\mathrm{Cov}(X_i, X_j) = -np_i p_j$ ($i \neq j$).
\end{tanim}

% =====================================================================
\section{Sürekli Dağılımlar}
\label{sec:surekli_dagilimlar}
% =====================================================================

\subsection{Düzgün (Uniform) Dağılım}
\label{subsec:uniform}

\begin{tanim}[Uniform Dağılım]
\label{def:uniform}
$X \sim \Uniform(a, b)$ ($a < b$):
\[
  f_X(x) = \frac{1}{b-a}, \quad a < x < b.
\]
\textbf{Beklenti:} $E[X] = (a+b)/2$. \quad
\textbf{Varyans:} $\mathrm{Var}(X) = (b-a)^2/12$. \quad
\textbf{MÜF:} $M_X(t) = \frac{e^{tb}-e^{ta}}{t(b-a)}$ ($t \neq 0$).
\end{tanim}

\begin{onerme}[Quantile Dönüşümü]
\label{prop:quantile}
$U \sim \Uniform(0,1)$ ve $F$ herhangi bir KDF ise $X = F^{-1}(U) \sim F$.
Bu simülasyonun temelidir: uniform rassal sayılardan herhangi bir dağılım üretilebilir.
\end{onerme}

\begin{proof}
$P(F^{-1}(U) \leq x) = P(U \leq F(x)) = F(x)$ ($U$ uniform olduğundan).
\end{proof}

\subsection{Normal Dağılım}
\label{subsec:normal}

\begin{tanim}[Normal (Gauss) Dağılımı]
\label{def:normal}
$X \sim \Normal(\mu, \sigma^2)$ ($\sigma > 0$):
\[
  f_X(x) = \frac{1}{\sigma\sqrt{2\pi}} \exp\!\left(-\frac{(x-\mu)^2}{2\sigma^2}\right), \quad x \in \mathbb{R}.
\]
\textbf{Beklenti:} $E[X] = \mu$. \quad
\textbf{Varyans:} $\mathrm{Var}(X) = \sigma^2$. \quad
\textbf{MÜF:} $M_X(t) = e^{\mu t + \sigma^2 t^2/2}$.\\
$Z = (X-\mu)/\sigma \sim \Normal(0,1)$ \textbf{standart normal}.
\end{tanim}

\begin{onerme}[Standart Normal Özellikleri]
$Z \sim \Normal(0,1)$. PDF: $\phi(z) = (2\pi)^{-1/2} e^{-z^2/2}$. KDF: $\Phi(z) = \int_{-\infty}^z \phi(t)\,dt$.
\begin{itemize}
  \item $\phi(-z) = \phi(z)$ (simetri); $\Phi(-z) = 1 - \Phi(z)$.
  \item $E[Z^{2k}] = (2k-1)!! = 1 \cdot 3 \cdots (2k-1)$; $E[Z^{2k+1}] = 0$.
  \item $E[Z^4] = 3$: basıklık 0 (referans).
\end{itemize}
\end{onerme}

\begin{teorem}[Normal Yeniden Üretkenliği]
\label{thm:normal_yeniden}
$X_i \sim \Normal(\mu_i, \sigma_i^2)$ bağımsız ($i=1,\ldots,n$) $\Rightarrow$
$\sum_{i=1}^n a_i X_i \sim \Normal\!\left(\sum a_i\mu_i,\ \sum a_i^2\sigma_i^2\right)$.
\end{teorem}

\begin{proof}
MÜF çarpım: $M_{\sum a_i X_i}(t) = \prod_i e^{\mu_i(a_i t) + \sigma_i^2(a_i t)^2/2}
= \exp\!\bigl(\sum a_i\mu_i \cdot t + \tfrac{t^2}{2}\sum a_i^2\sigma_i^2\bigr)$.
\end{proof}

\begin{tanim}[Çok Değişkenli Normal Dağılım]
\label{def:cok_normal}
$\mathbf{X} = (X_1, \ldots, X_n)^\top \sim \Normal_n(\boldsymbol{\mu}, \boldsymbol{\Sigma})$:
$\boldsymbol{\mu} \in \mathbb{R}^n$, $\boldsymbol{\Sigma}$ pozitif yarı-tanımlı $n\times n$ matris.
$\boldsymbol{\Sigma}$ pozitif tanımlı ise ortak PDF:
\[
  f(\mathbf{x}) = \frac{1}{(2\pi)^{n/2}|\boldsymbol{\Sigma}|^{1/2}}
  \exp\!\left(-\tfrac{1}{2}(\mathbf{x}-\boldsymbol{\mu})^\top \boldsymbol{\Sigma}^{-1}(\mathbf{x}-\boldsymbol{\mu})\right).
\]
$E[\mathbf{X}] = \boldsymbol{\mu}$, $\mathrm{Cov}(\mathbf{X}) = \boldsymbol{\Sigma}$.
\end{tanim}

\begin{teorem}[Çok Değişkenli Normal: Bağımsızlık]
\label{thm:cok_normal_bagimsiz}
$\mathbf{X} \sim \Normal_n(\boldsymbol{\mu}, \boldsymbol{\Sigma})$ ise
bileşenler çiftler arası korelasyonsuzluğu $\Leftrightarrow$ bağımsızlıktır:
$\Sigma_{ij} = 0$ ($i \neq j$) $\Leftrightarrow$ $X_i \perp X_j$.
\end{teorem}

\begin{proof}
$\boldsymbol{\Sigma}$ köşegen ise ortak PDF $\prod_i f_{X_i}(x_i)$ çarpımına ayrışır.
\end{proof}

\begin{hataanalizi}
Bu teorem genel dağılımlar için yanlıştır! Yalnızca çok değişkenli normal için
korelasyonsuzluk ile bağımsızlık eşdeğerdir. Genel dağılımlarda korelasyonsuz
ama bağımlı rasgele değişkenler inşa etmek mümkündür (bkz.\ Bölüm~\ref{sec:kovaryans}).
\end{hataanalizi}

\subsection{Üstel ve Gama Dağılımı}
\label{subsec:gama}

\begin{tanim}[Üstel Dağılım]
\label{def:exp}
$X \sim \Exp(\lambda)$ ($\lambda > 0$):
\[
  f_X(x) = \lambda e^{-\lambda x}, \quad x > 0.
\]
\textbf{Beklenti:} $E[X] = 1/\lambda$. \quad
\textbf{Varyans:} $\mathrm{Var}(X) = 1/\lambda^2$. \quad
\textbf{MÜF:} $M_X(t) = \lambda/(\lambda - t)$, $t < \lambda$.
\end{tanim}

\begin{onerme}[Bellek Yoksunluğu — Sürekli]
$X \sim \Exp(\lambda)$ için her $s, t > 0$'da $P(X > s+t \mid X > s) = P(X > t)$.
Sürekli dağılımlar arasında bellek yoksunluğu özelliğine sahip tek dağılım üsteldir.
\end{onerme}

\begin{proof}
$P(X > t) = e^{-\lambda t}$.
$P(X > s+t \mid X > s) = e^{-\lambda(s+t)}/e^{-\lambda s} = e^{-\lambda t} = P(X > t)$.
Tersine: $P(X > s+t) = P(X > s)P(X > t)$ ise $G(t) := P(X>t)$ için
$G(s+t) = G(s)G(t)$ ve süreklilik $G(t) = e^{-\lambda t}$ verir.
\end{proof}

\begin{tanim}[Gama Dağılımı]
\label{def:gama}
$X \sim \Gama(\alpha, \beta)$ ($\alpha, \beta > 0$):
\[
  f_X(x) = \frac{\beta^\alpha}{\Gamma(\alpha)} x^{\alpha-1} e^{-\beta x}, \quad x > 0.
\]
$\Gamma(\alpha) = \int_0^\infty t^{\alpha-1} e^{-t} \, dt$ (Gama fonksiyonu).
\textbf{Beklenti:} $E[X] = \alpha/\beta$. \quad
\textbf{Varyans:} $\mathrm{Var}(X) = \alpha/\beta^2$. \quad
\textbf{MÜF:} $M_X(t) = (\beta/(\beta-t))^\alpha$, $t < \beta$.\\
Özel durumlar: $\Gama(1, \beta) = \Exp(\beta)$; $\Gama(n/2, 1/2) = \chi^2_n$ (aşağıda).
\end{tanim}

\begin{proof}[Beklenti]
$E[X] = \frac{\beta^\alpha}{\Gamma(\alpha)} \int_0^\infty x^\alpha e^{-\beta x} dx
= \frac{\beta^\alpha}{\Gamma(\alpha)} \cdot \frac{\Gamma(\alpha+1)}{\beta^{\alpha+1}}
= \frac{\alpha}{\beta}$.
($\Gamma(\alpha+1) = \alpha\Gamma(\alpha)$ yinelemeli bağıntısı kullanıldı.)
\end{proof}

\begin{teorem}[Gama Yeniden Üretkenliği]
$X_i \sim \Gama(\alpha_i, \beta)$ bağımsız ($\beta$ aynı) $\Rightarrow$
$\sum X_i \sim \Gama(\sum \alpha_i, \beta)$.
\end{teorem}

\begin{proof}
$M_{\sum X_i}(t) = \prod_i (\beta/(\beta-t))^{\alpha_i} = (\beta/(\beta-t))^{\sum\alpha_i}$.
\end{proof}

\subsection{Beta Dağılımı}
\label{subsec:beta}

\begin{tanim}[Beta Dağılımı]
\label{def:beta}
$X \sim \Beta(\alpha, \beta)$ ($\alpha, \beta > 0$):
\[
  f_X(x) = \frac{x^{\alpha-1}(1-x)^{\beta-1}}{B(\alpha,\beta)}, \quad 0 < x < 1.
\]
$B(\alpha, \beta) = \Gamma(\alpha)\Gamma(\beta)/\Gamma(\alpha+\beta)$ (Beta fonksiyonu).\\
\textbf{Beklenti:} $E[X] = \alpha/(\alpha+\beta)$.\\
\textbf{Varyans:} $\mathrm{Var}(X) = \alpha\beta / [(\alpha+\beta)^2(\alpha+\beta+1)]$.\\
Özel durum: $\Beta(1,1) = \Uniform(0,1)$.
\end{tanim}

\begin{onerme}[Beta'nın Gama ile Bağlantısı]
$X \sim \Gama(\alpha, 1)$, $Y \sim \Gama(\beta, 1)$ bağımsız $\Rightarrow$
$U = X/(X+Y) \sim \Beta(\alpha, \beta)$, bağımsız olarak $X+Y \sim \Gama(\alpha+\beta, 1)$.
\end{onerme}

\begin{proof}
$(U, V) = (X/(X+Y), X+Y)$ dönüşümü; Jacobian $= v$;
$f_{U,V}(u,v) = f_{X,Y}(uv, v(1-u)) \cdot v$ hesaplanır;
$U$ ve $V$'nin bağımsız Beta ve Gama olduğu görülür.
\end{proof}

\subsection{Cauchy Dağılımı}
\label{subsec:cauchy}

\begin{tanim}[Cauchy Dağılımı]
\label{def:cauchy}
$X \sim \mathrm{Cauchy}(\mu, \sigma)$ ($\sigma > 0$):
\[
  f_X(x) = \frac{1}{\pi\sigma\left[1 + \left(\frac{x-\mu}{\sigma}\right)^2\right]}, \quad x \in \mathbb{R}.
\]
\textbf{Beklenti mevcut değildir} ($\int |x| f(x) dx = \infty$).\\
\textbf{KF:} $\varphi_X(t) = e^{i\mu t - \sigma|t|}$ (Laplace çekirdeği).\\
$\mathrm{Cauchy}(0,1)$: $Z_1/Z_2$ ($Z_1, Z_2$ bağımsız standart normal).
\end{tanim}

\begin{karsiornek}[Cauchy — Büyük Sayılar Kanunu Çöküyor]
$X_1, X_2, \ldots$ bağımsız $\mathrm{Cauchy}(0,1)$ ise
$\bar{X}_n = n^{-1}\sum X_i$ de $\mathrm{Cauchy}(0,1)$ dağılımlıdır —
ortalama $n$ büüdükçe stabilize olmaz. Büyük Sayılar Kanunu beklentinin
varlığını gerektirir; Cauchy bunu ihlal eder.
\end{karsiornek}

\subsection{Lognormal Dağılım}
\label{subsec:lognormal}

\begin{tanim}[Lognormal Dağılım]
\label{def:lognormal}
$Y = e^X$, $X \sim \Normal(\mu, \sigma^2)$ ise $Y \sim \mathrm{Lognormal}(\mu, \sigma^2)$:
\[
  f_Y(y) = \frac{1}{y\sigma\sqrt{2\pi}} \exp\!\left(-\frac{(\ln y - \mu)^2}{2\sigma^2}\right), \quad y > 0.
\]
\textbf{Beklenti:} $E[Y] = e^{\mu + \sigma^2/2}$.\\
\textbf{Varyans:} $\mathrm{Var}(Y) = e^{2\mu+\sigma^2}(e^{\sigma^2}-1)$.\\
Tüm momentler mevcuttur ama MÜF yoktur ($E[e^{tY}] = \infty$ her $t > 0$ için).
\end{tanim}

\begin{ornek}[Finans Uygulaması]
Black-Scholes modelinde hisse senedi fiyatı $S_t = S_0 e^{(\mu-\sigma^2/2)t + \sigma W_t}$
lognormal dağılım izler. $\ln(S_t/S_0) \sim \Normal((\mu-\sigma^2/2)t, \sigma^2 t)$.
\end{ornek}

\subsection{Chi-Kare, \texorpdfstring{$t$}{t} ve \texorpdfstring{$F$}{F} Dağılımları}
\label{subsec:chi_t_f}

\begin{tanim}[Chi-Kare Dağılımı]
\label{def:chi_kare}
$Z_1, \ldots, Z_n$ bağımsız $\Normal(0,1)$ ise
\[
  \chi^2_n = \sum_{i=1}^n Z_i^2 \sim \Gama(n/2,\, 1/2).
\]
\textbf{Beklenti:} $E[\chi^2_n] = n$. \quad
\textbf{Varyans:} $\mathrm{Var}(\chi^2_n) = 2n$. \quad
\textbf{MÜF:} $M(t) = (1-2t)^{-n/2}$, $t < 1/2$.
\end{tanim}

\begin{proof}
$Z_i^2 \sim \Gama(1/2, 1/2)$ (Bölüm~\ref{sec:rd_donusum}'deki $Y=X^2$ dönüşümü).
Gama yeniden üretkenliği: $\sum Z_i^2 \sim \Gama(n/2, 1/2)$.
\end{proof}

\begin{tanim}[$t$ Dağılımı (Student)]
\label{def:t_dagilim}
$Z \sim \Normal(0,1)$, $V \sim \chi^2_n$ bağımsız. $T = Z/\sqrt{V/n} \sim t_n$ (Student-$t$).
PDF:
\[
  f_T(t) = \frac{\Gamma((n+1)/2)}{\sqrt{n\pi}\,\Gamma(n/2)}\left(1+\frac{t^2}{n}\right)^{-(n+1)/2}, \quad t \in \mathbb{R}.
\]
$n = 1$: $t_1 = \mathrm{Cauchy}(0,1)$.
$n \to \infty$: $t_n \xrightarrow{d} \Normal(0,1)$ (Böl.~\ref{chp:merkezi_limit}).\\
\textbf{Beklenti:} $E[T] = 0$ ($n > 1$). \quad \textbf{Varyans:} $\mathrm{Var}(T) = n/(n-2)$ ($n > 2$).
\end{tanim}

\begin{tanim}[$F$ Dağılımı]
\label{def:f_dagilim}
$U \sim \chi^2_m$, $V \sim \chi^2_n$ bağımsız.
$W = (U/m)/(V/n) \sim F(m, n)$.\\
\textbf{Beklenti:} $E[W] = n/(n-2)$ ($n > 2$).\\
Bağlantı: $T_n^2 \sim F(1, n)$.
\end{tanim}

\begin{hocanotu}
$t_n$ ve $F(m,n)$ dağılımları istatistiksel hipotez testinin (Bölüm~\ref{chp:hipotez_testi})
temel test istatistikleridir. Normal popülasyondan çekilen örneklem ortalaması
$\bar{X}$ ve örneklem varyansı $S^2$ ile bunlar türetilir (Böl.~\ref{chp:istatistik_temelleri}).
\end{hocanotu}

% =====================================================================
\section{Üstel Aile}
\label{sec:ustel_aile}
% =====================================================================

\begin{kesif}
Binom, Poisson, Normal, Gama, Beta — bu dağılımların PDF/PMF'leri çok farklı görünüyor.
Ama hepsinin $\exp(\ldots)$ biçiminde yazılabildiğini fark ettiniz mi?
Bu yapı tesadüf mü, yoksa derin bir ortak özelliği mi yansıtıyor?
\end{kesif}

\begin{tanim}[Tek Parametreli Üstel Aile]
\label{def:ustel_aile}
$\theta$'ya bağlı bir dağılım ailesi, eğer PDF/PMF'si
\[
  f(x \mid \theta) = h(x) \exp\!\bigl(\eta(\theta) T(x) - B(\theta)\bigr)
\]
biçiminde yazılabiliyorsa \textbf{(tek parametreli) üstel aile} (exponential family) üyesidir.
Burada:
\begin{itemize}
  \item $\eta(\theta)$: \textbf{doğal parametre} (natural parameter),
  \item $T(x)$: \textbf{yeterli istatistik} (sufficient statistic),
  \item $B(\theta)$: normalleştirme sabiti (log-bölüşüm fonksiyonu),
  \item $h(x)$: $\theta$'dan bağımsız taban ölçüsü.
\end{itemize}
\textbf{Çok parametreli genelleştirme:}
$f(x \mid \boldsymbol{\theta}) = h(x) \exp\!\bigl(\sum_{j=1}^k \eta_j(\boldsymbol{\theta}) T_j(x) - B(\boldsymbol{\theta})\bigr)$.
\end{tanim}

\begin{ornek}[Üstel Aile Örnekleri]
\begin{itemize}
  \item \textbf{Normal} ($\sigma^2$ biliniyor): $\eta(\mu) = \mu/\sigma^2$, $T(x) = x$, $B(\mu) = \mu^2/(2\sigma^2)$.
  \item \textbf{Poisson}: $\eta(\lambda) = \ln\lambda$, $T(k) = k$, $B(\lambda) = \lambda$, $h(k) = 1/k!$.
  \item \textbf{Binom} ($n$ sabit): $\eta(p) = \ln(p/q)$ (log-odds), $T(k) = k$.
  \item \textbf{Gama} ($\beta$ biliniyor): $\eta(\alpha) = \alpha-1$, $T(x) = \ln x$.
\end{itemize}
\end{ornek}

\begin{teorem}[Log-Bölüşüm Fonksiyonunun Türevleri]
\label{thm:log_bolüşüm}
Doğal parametre formu $f(x \mid \eta) = h(x)\exp(\eta T(x) - B(\eta))$ için:
\[
  E_\eta[T(X)] = B'(\eta), \qquad \mathrm{Var}_\eta(T(X)) = B''(\eta).
\]
\end{teorem}

\begin{proof}
$1 = \int h(x) e^{\eta T(x) - B(\eta)} dx$ her $\eta$ için.
$\eta$'ya göre türev: $0 = \int T(x) h(x) e^{\eta T - B} dx - B'(\eta) \cdot 1$,
yani $B'(\eta) = E[T(X)]$.
İkinci türev: $B''(\eta) = E[T(X)^2] - (E[T(X)])^2 = \mathrm{Var}(T(X))$.
\end{proof}

\begin{onerme}[Üstel Ailenin İstatistik Teorisindeki Önemi]
Üstel aile üyesi dağılımlar için $\sum_{i=1}^n T(X_i)$ yeterli istatistiktir
(Neyman faktorizasyon teoremi, Böl.~\ref{chp:nokta_tahmin}).
Bu dağılımlar UMVUE tahminlerine (Rao-Blackwell, Böl.~\ref{chp:nokta_tahmin}) sahiptir.
\end{onerme}

% =====================================================================
\section{Yer-Ölçek Aileleri}
\label{sec:yer_olcek}
% =====================================================================

\begin{tanim}[Yer-Ölçek Ailesi]
\label{def:yer_olcek}
$f_0$ bir "standart" PDF olsun. \textbf{Yer-ölçek ailesi}:
\[
  f(x \mid \mu, \sigma) = \frac{1}{\sigma} f_0\!\left(\frac{x-\mu}{\sigma}\right),
  \quad \mu \in \mathbb{R},\ \sigma > 0.
\]
$\mu$: yer parametresi (konum, location); $\sigma$: ölçek parametresi (scale).
\end{tanim}

\begin{ornek}[Yer-Ölçek Aileleri]
\begin{itemize}
  \item Normal: $f_0(z) = \phi(z)$; $\Normal(\mu, \sigma^2)$ yer-ölçek ailesi.
  \item Üstel (ölçek): $f_0(z) = e^{-z}$; $\Exp(\lambda) = \mathrm{Ölçek}(1/\lambda)$ ailesi.
  \item Cauchy: $f_0(z) = [\pi(1+z^2)]^{-1}$; $\mathrm{Cauchy}(\mu,\sigma)$ yer-ölçek.
  \item Logistik: $f_0(z) = e^{-z}/(1+e^{-z})^2$; lojistik regresyon temeli.
\end{itemize}
\end{ornek}

\begin{onerme}[Standartlaştırma]
$(X - \mu)/\sigma$ "standart" dağılıma sahip; yani yer-ölçek ailelerinde
tek bir "standart" tablo yeterlidir: $F(x \mid \mu, \sigma) = F_0((x-\mu)/\sigma)$.
\end{onerme}

% ---- Disiplinlerarası Bağlantı (TYMM) ----
\begin{kopru}[Disiplinlerarası — Biyoloji, Fizik, Mühendislik]
\textbf{Biyoloji:} Popülasyon büyüklüğü, hücre bölünme süreleri üstel ve gama
dağılımlarıyla modellenir; lognormal dağılım protein konsantrasyonlarını açıklar.\\
\textbf{Fizik:} Maxwell-Boltzmann hız dağılımı gama ailesine aittir;
termal gürültü normal dağılımdan kaynaklanır (Merkezi Limit Teoremi).\\
\textbf{Mühendislik/Güvenilirlik:} Weibull dağılımı (Gama genellemesi)
ürün ömrü modellemesinin standardıdır; üstel dağılım "bellek yoksunluğu"
sayesinde sıra kuyruk modellerinin temelidir.
\end{kopru}

% ---- Öz-Yansıtma ----
\begin{ozyansima}
\begin{itemize}[leftmargin=1.8em]
  \item Binom ile Poisson arasındaki "limit" ilişkisini kendi cümlelerinizle açıklayın.
        Hangi koşullar altında Poisson yaklaşımı makul bir seçimdir?
  \item Bellek yoksunluğu özelliği fiziksel olarak ne anlama gelir?
        Hangi gerçek hayat fenomenlerini modelleyebilir, hangilerini modelleyemez?
  \item Üstel aile kavramı neden istatistik teorisi açısından önemlidir?
  \item $t$ ve $F$ dağılımları neden "örneklem dağılımları" olarak anılır?
\end{itemize}
\end{ozyansima}

% ---- Köprü Notu ----
\begin{kopru}
\textbf{Önceki bölüm:} Bölüm~\ref{chp:beklenti} — beklenti, MÜF, KF araçları
buradaki her dağılım için somutlaştırıldı.\\
\textbf{Sonraki bölüm:} Bölüm~\ref{chp:buyuk_sayilar} — büyük sayılar kanunlarında
bu dağılımların asimptotik davranışı; Bölüm~\ref{chp:merkezi_limit}'te normal
dağılıma yakınsama.\\
\textbf{İstatistik kısmı:} Böl.~\ref{chp:nokta_tahmin}'de üstel aile ve yeterli
istatistik; Böl.~\ref{chp:hipotez_testi}'nde $\chi^2$, $t$, $F$ test istatistikleri.
\defterbaglantisi{bolum05\_dagilimlar.ipynb}{%
  Dağılım grafikleri, parametre duyarlılığı, simülasyon karşılaştırması}
\end{kopru}

% =====================================================================
%  ALIŞTIRMALAR
% =====================================================================

\cozumlualistirmalar

\begin{alistirma}[Al.5.1]{A}{Poisson Hesabı}
Birim saatte ortalama 5 arıza oluşan bir sistemde:
(a) Tam 3 arıza olasılığı. (b) En az 1 arıza olasılığı. (c) En fazla 2 arıza olasılığı.
\end{alistirma}
\alcozum{%
$X \sim \Pois(5)$.\\
(a) $P(X=3) = e^{-5}5^3/3! = 125e^{-5}/6 \approx 0.1404$.\\
(b) $P(X \geq 1) = 1 - P(X=0) = 1 - e^{-5} \approx 0.9933$.\\
(c) $P(X \leq 2) = e^{-5}(1 + 5 + 25/2) = 18.5e^{-5} \approx 0.1247$.}

\begin{alistirma}[Al.5.2]{A}{Normal Dağılım}
$X \sim \Normal(70, 100)$ (not: sınav puanı, $\sigma = 10$).
(a) $P(X > 85)$. (b) $P(60 < X < 80)$. (c) 90. yüzdeliği bulun.
\end{alistirma}
\alcozum{%
$Z = (X-70)/10 \sim \Normal(0,1)$.\\
(a) $P(X>85) = P(Z > 1.5) = 1 - \Phi(1.5) \approx 1 - 0.9332 = 0.0668$.\\
(b) $P(60<X<80) = P(-1 < Z < 1) = 2\Phi(1) - 1 \approx 2(0.8413) - 1 = 0.6827$.\\
(c) $\Phi(z_{0.9}) = 0.9 \Rightarrow z_{0.9} \approx 1.282$; $x_{0.9} = 70 + 10(1.282) = 82.82$.}

\begin{alistirma}[Al.5.3]{B}{Gama Dağılımı ve Beklenti}
$X \sim \Gama(3, 2)$.
(a) $E[X]$, $\mathrm{Var}(X)$ ve $M_X(t)$'yi yazın.
(b) $P(X > 3)$'ü sayısal ifade edin (yarı analitik form yeterli).
(c) $X_1, X_2 \sim \Gama(3, 2)$ bağımsız ise $X_1 + X_2$'nin dağılımını bulun.
\end{alistirma}
\alcozum{%
(a) $E[X] = 3/2$, $\mathrm{Var}(X) = 3/4$, $M_X(t) = (2/(2-t))^3$ ($t < 2$).\\
(b) $P(X>3) = \int_3^\infty \frac{2^3}{\Gamma(3)} x^2 e^{-2x} dx = 1 - F_{\Gama(3,2)}(3)$
(tamamlanmamış gama fonksiyonu).\\
(c) $X_1+X_2 \sim \Gama(6, 2)$.}

\begin{alistirma}[Al.5.4]{B}{Üstel Ailede Log-Bölüşüm}
$X \sim \Pois(\lambda)$'yı doğal parametre $\eta = \ln\lambda$ ile üstel aile formunda yazın.
$B(\eta)$'yı bulun ve $B'(\eta) = E[T(X)]$ formülüyle $E[X] = \lambda$'yı doğrulayın.
\end{alistirma}
\alcozum{%
$P(X=k) = \frac{\lambda^k e^{-\lambda}}{k!} = \frac{1}{k!} \exp(k\ln\lambda - \lambda)
= \frac{1}{k!} \exp(\eta k - e^\eta)$.\\
$T(k) = k$, $h(k) = 1/k!$, $B(\eta) = e^\eta = \lambda$.\\
$B'(\eta) = e^\eta = \lambda = E[T(X)] = E[X]$. \checkmark}

% ---

\alistirmalar

\begin{alistirma}[Al.5.5]{A}{Binom Olasılık}
10 soru içeren çoktan seçmeli testte her sorunun 4 şıkkı var; öğrenci tamamen rastgele cevaplıyor.
(a) Tam 3 doğru cevap olasılığı. (b) En az 5 doğru cevap olasılığı.
(c) $E[X]$ ve $\mathrm{Var}(X)$.
\end{alistirma}
\alcozum{%
$X \sim \Binom(10, 1/4)$.
(a) $\binom{10}{3}(1/4)^3(3/4)^7 \approx 0.2503$.
(b) $P(X\geq5) = \sum_{k=5}^{10}\binom{10}{k}(1/4)^k(3/4)^{10-k} \approx 0.0781$.
(c) $E[X] = 2.5$, $\mathrm{Var}(X) = 1.875$.}

\begin{alistirma}[Al.5.6]{A}{Geometrik Uygulama}
Bir üretim hattında her ürün $p = 0.05$ olasılıkla bozuk.
(a) İlk bozuk ürünün 10. üründe çıkma olasılığı.
(b) İlk bozuk üründen önce en az 20 ürün üretilme olasılığı.
(c) Kaçıncı üründe ilk bozuğun çıkması beklenir?
\end{alistirma}
\alcozum{%
$X \sim \mathrm{Geom}(0.05)$.
(a) $P(X=10) = (0.95)^9(0.05) \approx 0.0315$.
(b) $P(X>20) = (0.95)^{20} \approx 0.3585$.
(c) $E[X] = 1/0.05 = 20$.}

\begin{alistirma}[Al.5.7]{B}{Beta Dağılımı}
$X \sim \Beta(2, 3)$.
(a) $E[X]$ ve $\mathrm{Var}(X)$'i hesaplayın.
(b) $P(X > 0.5)$'i analitik olarak bulun.
(c) Bu dağılım hangi pratik durumu modelleyebilir?
\end{alistirma}
\alcozum{%
(a) $E[X] = 2/5 = 0.4$. $\mathrm{Var}(X) = 2\cdot3/(25\cdot6) = 6/150 = 1/25$.\\
(b) $f(x) = 12x(1-x)^2$.
$P(X>0.5) = 12\int_{0.5}^1 x(1-x)^2 dx = 12[x^2/2 - 2x^3/3 + x^4/4]_{0.5}^1
= 12[(1/2-2/3+1/4)-(1/8-1/12+1/16)] = 5/16$.\\
(c) $[0,1]$'de değer alan oranlar: dönüşüm hata oranı, bir proje tamamlanma yüzdesi, vb.}

\begin{alistirma}[Al.5.8]{B}{$\chi^2$ ve $t$ Dağılımı}
$X_1, \ldots, X_5$ bağımsız $\Normal(0,1)$.
(a) $V = \sum X_i^2$'nin dağılımını yazın; $E[V]$ ve $\mathrm{Var}(V)$'yi bulun.
(b) $Z \sim \Normal(0,1)$ bu $X_i$'lerden bağımsız. $T = Z/\sqrt{V/5}$'in dağılımını yazın.
(c) $T^2$'nin dağılımını yazın.
\end{alistirma}
\alcozum{%
(a) $V \sim \chi^2_5 = \Gama(5/2, 1/2)$. $E[V] = 5$, $\mathrm{Var}(V) = 10$.\\
(b) $T \sim t_5$.\\
(c) $T^2 \sim F(1, 5)$.}

\begin{alistirma}[Al.5.9]{B}{Poisson-Gama Karışımı}
$\Lambda \sim \Gama(\alpha, \beta)$ ve $X \mid \Lambda = \lambda \sim \Pois(\lambda)$.
$X$'in koşulsuz dağılımının negatif binom olduğunu gösterin.
\end{alistirma}
\alcozum{%
$P(X=k) = \int_0^\infty \frac{\lambda^k e^{-\lambda}}{k!} \cdot \frac{\beta^\alpha}{\Gamma(\alpha)}\lambda^{\alpha-1}e^{-\beta\lambda} d\lambda
= \frac{\beta^\alpha}{k!\,\Gamma(\alpha)} \int_0^\infty \lambda^{k+\alpha-1} e^{-(\beta+1)\lambda} d\lambda
= \frac{\beta^\alpha}{k!\,\Gamma(\alpha)} \cdot \frac{\Gamma(k+\alpha)}{(\beta+1)^{k+\alpha}}
= \binom{k+\alpha-1}{k} \left(\frac{\beta}{\beta+1}\right)^\alpha \left(\frac{1}{\beta+1}\right)^k$.
Bu $\mathrm{NegBin}(\alpha, \beta/(\beta+1))$ dağılımıdır.}

\begin{alistirma}[Al.5.10]{B}{Yer-Ölçek Dönüşümü}
$X \sim \mathrm{Cauchy}(0,1)$ ise $Y = \mu + \sigma X$ ($\sigma > 0$) için
$Y \sim \mathrm{Cauchy}(\mu, \sigma)$ olduğunu gösterin.
Cauchy dağılımı neden bir yer-ölçek ailesidir?
\end{alistirma}
\alcozum{%
$f_X(x) = 1/(\pi(1+x^2))$. Jacobian dönüşümü: $x = (y-\mu)/\sigma$, $|dx/dy| = 1/\sigma$.
$f_Y(y) = \frac{1}{\pi(1+((y-\mu)/\sigma)^2)} \cdot \frac{1}{\sigma}
= \frac{1}{\pi\sigma(1+((y-\mu)/\sigma)^2)}$. Bu $\mathrm{Cauchy}(\mu,\sigma)$ PDF'sidir.
Yer-ölçek ailesi: standart form $f_0(z) = 1/(\pi(1+z^2))$; $f(y|\mu,\sigma) = f_0((y-\mu)/\sigma)/\sigma$.}

\begin{alistirma}[Al.5.11]{C}{Normal Dağılımın Üstel Aile Formu}
$X \sim \Normal(\mu, \sigma^2)$'yi iki parametreli üstel ailede yazın:
$\boldsymbol{\eta} = (\eta_1, \eta_2)$, $\mathbf{T}(x) = (T_1(x), T_2(x))$ ile.
$B(\boldsymbol{\eta})$'yı bulun ve $\partial B/\partial\eta_1 = E[T_1(X)]$,
$\partial B/\partial\eta_2 = E[T_2(X)]$ formüllerini doğrulayın.
\end{alistirma}
\alcozum{%
$f(x|\mu,\sigma^2) = \frac{1}{\sqrt{2\pi}\sigma}\exp\!\left(-\frac{x^2}{2\sigma^2} + \frac{\mu}{\sigma^2}x - \frac{\mu^2}{2\sigma^2}\right)$.
Doğal parametreler: $\eta_1 = \mu/\sigma^2$, $\eta_2 = -1/(2\sigma^2)$ (yani $\sigma^2 = -1/(2\eta_2)$, $\mu = \eta_1/(-2\eta_2)$).
$T_1(x) = x$, $T_2(x) = x^2$.
$B(\boldsymbol{\eta}) = -\eta_1^2/(4\eta_2) - \frac{1}{2}\ln(-2\eta_2) + \frac{1}{2}\ln(2\pi)$.
$\partial B/\partial\eta_1 = -\eta_1/(2\eta_2) = \mu/\sigma^2 \cdot \sigma^2 = \mu = E[X]$. \checkmark
$\partial B/\partial\eta_2 = \eta_1^2/(4\eta_2^2) - 1/(2\eta_2) = \mu^2 + \sigma^2 = E[X^2]$. \checkmark}

\begin{alistirma}[Al.5.12]{C}{Binom $\to$ Poisson: Tam İspat}
Teorem~\ref{thm:binom_poisson}'in ispatında kullanılan iki limiti kesin biçimde ispatlayın:
(a) $n(n-1)\cdots(n-k+1)/n^k \to 1$.
(b) $(1 - \lambda_n/n)^n \to e^{-\lambda}$ ($\lambda_n \to \lambda$).
\end{alistirma}
\alcozum{%
(a) $\frac{n(n-1)\cdots(n-k+1)}{n^k} = 1 \cdot (1-1/n) \cdots (1-(k-1)/n) \to 1^k = 1$ (sabit $k$, $n\to\infty$).\\
(b) $\ln(1-\lambda_n/n)^n = n\ln(1-\lambda_n/n)$. Taylor: $\ln(1-x) = -x - x^2/2 - \cdots$,
$n \cdot (-\lambda_n/n + O(1/n^2)) = -\lambda_n + O(1/n) \to -\lambda$.
Dolayısıyla $(1-\lambda_n/n)^n \to e^{-\lambda}$.}

% =====================================================================
\endinput
